\documentclass[11pt]{article}
\usepackage[T1]{fontenc}
\usepackage{amsmath,amssymb,amsthm,booktabs,geometry,array,url}
\geometry{margin=1in}
\usepackage[hidelinks]{hyperref}
\hypersetup{pdftitle={Irrationality in the five odd zeta values from 7 to 15},
  pdfauthor={JanaHesse580},pdfsubject={B40a determinant construction and Lean 4 formalization}}
\newtheorem{theorem}{Theorem}
\newtheorem{proposition}[theorem]{Proposition}
\newcommand{\Z}{\mathbb Z}
\newcommand{\Q}{\mathbb Q}
\newcommand{\R}{\mathbb R}
\newcommand{\C}{\mathbb C}
\newcommand{\val}{v_p}
\newcommand{\dd}{\,\mathrm d}
\newcommand{\code}[1]{\texttt{\detokenize{#1}}}
\setlength{\emergencystretch}{2em}
\title{Irrationality in the five odd zeta values from 7 to 15\\
\large A determinant construction with a Lean 4 certificate}
\author{JanaHesse580\\ \small\texttt{JanaHesse580@gmail.com}}
\date{November 2026 \qquad Software release 1.0.0}
\begin{document}
\maketitle
\begin{abstract}
For the Riemann zeta function, at least one of the values at
$7,9,11,13,15$ is irrational. We use a determinant of order $40n$ whose
entries are rational linear forms in these values and $1$. At the scale
$2000n^2$ the analytic determinant has upper exponent $-6.19$, while
clearing its polynomial content costs at most $6.1405$. A confluent
interpolation calculation at primes $51n+1$ supplies nonvanishing at
rational specializations. The accompanying Lean development formalizes
the argument. One finite check on 1,375 arithmetic rows uses native
evaluation. A separate standard-library Python program recomputes these
rows and the numerical bounds.
\end{abstract}

\section{Statement and the integer contradiction}
Zudilin's four-value result concerns $\{5,7,9,11\}$~\cite{zud};
his earlier theorem also gives an irrational value among the odd indices
from 7 to 37~\cite{zud2}. These statements illustrate the importance
of specifying both endpoints. The six-value Lean development used here
as a software baseline includes the index 17~\cite{six}. We change the
kernel, the size of the determinant, the local interpolation problem and
the bounds on the determinant and its content.

\begin{theorem}\label{thm}
There exists $s\in\{7,9,11,13,15\}$ for which $\zeta(s)$ is irrational.
\end{theorem}
At these real arguments the Dirichlet series converges to the real part
of Mathlib's complex Riemann zeta function. This is the quantity in the
formal statement.

Put $h=40n$ and $N=2000n^2$. The construction gives a rational polynomial
determinant $D_n(\mathbf X)=c_n P_n(\mathbf X)$, with $c_n\ne0$,
$P_n\in\Z[X_1,\ldots,X_5]$, and $\deg P_n\le h$. Write $\mathbf z$ for
the ordered vector of the five zeta values. The three assertions needed are
\begin{align}
 D_n(\mathbf z)&=\det H_n,\qquad
 |\det H_n|\le \exp(-6.19N) && (n\gg1),\label{eq:A}\\
 |c_n|^{-1}&\le \exp(6.1405N) && (n\gg1),\label{eq:C}\\
 P_n(\mathbf q)&\ne0 &&(n\text{ in a subsequence}).\label{eq:gate}
\end{align}
Suppose $Q\mathbf z\in\Z^5$ for an integer $Q\ge1$. Then
$Q^h P_n(\mathbf z)$ is an integer. By \eqref{eq:gate} it is nonzero
on a subsequence. Equations \eqref{eq:A}--\eqref{eq:C} bound its absolute
value above by $\exp(40n\log Q-0.0495N)$, which tends to zero.
This proves Theorem~\ref{thm} once the three assertions are established.

\section{The B40a kernel and the five coordinates}
For $(t)_m=\prod_{j=0}^{m-1}(t+j)$, let $d$ run with multiplicities
over $\mathcal D=(0^6,8^4,9^2)$, where the superscripts denote repetitions.
Set $c=59n+1$, $L=50$, $U=34$ and $s_0=25$. Define
\begin{align*}
 K_n(t)&=\frac{2\prod_{d\in\mathcal D}((50-2d)n)!}{((25n)!)^{10}}\,
 \frac{(t+c)\big[(t+1)_{25n}(t+93n+2)_{25n}\big]^5}
 {\prod_{d\in\mathcal D}(t+(34+d)n+1)_{(50-2d)n+1}},\\
 E_0(x)&=1,\qquad E_i(x)=\frac{2x^2}{(2i)!}
                    \prod_{j=1}^{i-1}(x^2-j^2)\quad(i\ge1),\\
 R_{ij}(t)&=K_n(t)E_i(t+c)E_j(t+c),\qquad
 (H_n)_{ij}=\sum_{v=0}^{\infty} R_{ij}^{(4)}(v-25n)/24,\quad 0\le i,j<h .
\end{align*}
The denominator has degree $500n+12$, the unweighted numerator degree
$250n+1$, and the basis factors together have degree at most $160n-4$ for
$n\ge1$. In particular the rational function decays faster than $1/t$.
Write its principal parts as $\sum_{k,\ell} a_{k,\ell}(t+k)^{-\ell}$.
Multiplication by the common denominator and comparison of the highest
possible coefficient give $\sum_k a_{k,1}=0$. Reflection
$t\mapsto-2c-t$ pairs $k$ with $2c-k$. The two numerator blocks exchange,
the centered linear factor changes sign and the even basis factors are
invariant. The reflected coefficients satisfy
$a_{2c-k,\ell}=(-1)^{\ell+1}a_{k,\ell}$. All aggregate even orders
therefore vanish. Every pole has order at most 12. The surviving orders
are consequently $3,5,7,9,11$.

After four derivatives the contribution of $(t+k)^{-\ell}$ is
$\binom{\ell+3}{4}(t+k)^{-\ell-4}$. Its sum is that coefficient times
$\zeta(\ell+4)-\sum_{r=1}^{k-25n-1}r^{-\ell-4}$.
Thus order one cancels the possible $\zeta(5)$, and order thirteen,
which would give $\zeta(17)$, is absent. Replacing the remaining five
zeta values by variables gives an affine polynomial matrix $H^{\rm pol}$.
Taking its determinant proves the evaluation identity and the degree
bound. Dividing by rational content gives the integer polynomial used
in Section~1. The formal construction also chooses a nonzero content
in the degenerate case of a zero polynomial.

\section{The local nonvanishing calculation}
Take $p=51n+1$ prime and $\mathbf q$ integral at $p$. Work first in
the basis $x^{2i}$. For $p>50n$ all factors from the within-block
factorials are $p$-units. Split each finite harmonic sum into its
$p$-integral terms and its denominators divisible by $p$. At this $p$,
the singular part is supported on $k=76n+2,\ldots,84n+1$.
Indeed $k-25n-1\ge p$ begins exactly at the first listed pole, while
$k-25n-1<2p$. These $8n$ poles have multiplicity six. At each one
the constant coefficient $u_k$ of the regularized kernel series is a
unit. Removing the vanishing factors from the denominator leaves
nonzero factors modulo $p$; the first numerator product at $-k$ has
absolute factors $k-1-j$ strictly between $p$ and $2p$
($0\le j<25n$). The second has factors $93n+2+j-k$ strictly between
zero and $p$. The centered linear factor also has absolute value
strictly between zero and $p$. This proves the claimed unit condition.

Index columns by a pole and a Taylor order $a=0,\dots,4$. The associated
change of basis is invertible modulo $p$. In absolute value the signed
nodes $c-k$ range from $17n+1$ to $25n$ and their squares are distinct
modulo $p$: the nodes are distinct and the sum of two absolute values
is at most $50n<p$. The derivative of the squaring map is a unit at
each node. If the first five coefficients of $F(x^2)$ vanish there,
then $(X-(c-k)^2)^5$ divides $F$. The product over $8n$ nodes has degree
$40n$, forcing $F=0$ when $\deg F<40n$. This proves invertibility of
the square interpolation matrix.

In interpolated coordinates scale each row and column of order $a$ by
$p^{5-a}$. Off-diagonal pole blocks vanish modulo $p$. At a single
pole the singular term of combined order $a+b$ leaves, up to the common
sign and the unit $u_k$, the integer matrix
\[
 B=\left(\binom{9-a-b}{4}\right)_{0\le a,b<5}
 =\begin{pmatrix}126&70&35&15&5\\70&35&15&5&1\\35&15&5&1&0\\
 15&5&1&0&0\\5&1&0&0&0\end{pmatrix}.
\]
Here the fourth derivative supplies the binomial coefficient; positive
powers of $p$ kill higher Taylor corrections. More explicitly, if
$U_k(T)$ is the quotient of the translated numerator by the translated
pole-free cofactor, truncated to the six local coefficients, the scaled
entry before reduction has weight
\[
 \sum_{r=0}^{5-a-b}p^r\binom{9-a-b-r}{4}[T^r]U_k(T)
\]
with an empty sum for $a+b>5$. All coefficients lie in the localization
at $p$ and $[T^0]U_k=u_k$ modulo $p$. The interpolating polynomials
have prescribed jets at this pole and zero jets through order four at
the others. Terms from other poles and from the regular harmonic part
drop out after the scaling and reduction (the regular part carries
$p^{10-a-b}$). To compute the determinant,
replace the rows by their successive forward differences evaluated at row
zero, and then do the same to the columns. Both transformations have
determinant one. Pascal's identity gives zero for $a+b>4$ and one for
$a+b=4$. Reversal of five columns has sign $(-1)^{10}=1$, so $\det B=1$.
The actual reduced block is $-u_k B$ and its determinant is a unit.
This distinction matters for the valuation.

Each diagonal scaling contributes $p^{(5+4+3+2+1)8n}=p^{120n}$.
After undoing both scalings the monomial determinant has valuation
$-240n$. Moving back to $E_i$ multiplies it by
$\prod_{i=1}^{40n-1}(2/(2i)!)^2$. Since $2i<2p$, exactly the indices
$i\ge(p+1)/2$ contribute one factor of $p$ in $(2i)!$.
Primality forces $n$ even, so there are $(29n-2)/2$ such indices.
We conclude
\[
 \val D_n(\mathbf q)=-240n-(29n-2)=-269n+2 .
\]
There are arbitrarily large primes congruent to $1\bmod 51$.
Choose them greater than a fixed common denominator $Q$ to ensure
$\mathbf q$ is integral at $p$. This proves \eqref{eq:gate}.

\section{Contour and energy bound}
The vertical residue integral starts at real part
$\sigma_n=-\lfloor17n/20\rfloor-1/2$. Quintuple zeros of the numerator
remove the unwanted initial fourth-derivative terms. Closing a rectangle
and passing to the height limit gives the moment matrix multiplied by
$-2\pi i$. Deform the two halves separately and put $D=50n$. In each half
the eight nodes have real part $\sigma_n+D(0.017-c_j)$ and imaginary
part $\pm D y_j$, with rational data
\[
\begin{array}{c|rrrrrrrr}
j&0&1&2&3&4&5&6&7\\ \hline
1000y_j&0&25&50&100&200&350&550&850\\
1000c_j&17&22&10&-20&-81&-163&-257&-455
\end{array}
\]
Join adjacent nodes by line segments and continue vertically from the last
one. Holomorphy in the deformation domains and decay on the closing
segments justify the replacement. On the truncated contour write
$z_j=((t_j+c)/D)^2$ and $\Delta(z)=\prod_{i<j}(z_j-z_i)$.
Expansion of the determinant, Fubini's theorem and symmetry of the
product measure give the moment identity
\[
 h!\det H_{n,T}^{\rm bend}=
 B_n\int\Delta(z)^2\prod_j w_n(t_j)\,\dd t_j,\qquad
 B_n=\prod_{i=0}^{h-1}\big(\operatorname{lc}(E_i)D^{2i}\big)^2 .
\]
The scalar weight on each oriented half is
\[
 w_n(t)\dd t=K_n(t)\,\frac{1}{24}
 \frac{\mathrm d^4}{\mathrm dt^4}\big(\pi\cot(\pi t)\big)\dd t .
\] The limit on the left, after division by $h!$, is
$(-2\pi i)^h\det H_n$.

The certificate supplies 34 triples of center, radius and mass in one
half-plane. Reflect them, giving each copy mass $m_i/2000$ from the
integer mass $m_i$; the upper masses sum to 1000. Spread each mass
uniformly over its circle and denote the probability measure by $\mu$.
Its potential and energy in the positive-log convention are
\[
 U(z)=\sum b_i\log\max(r_i,|z-u_i|),\qquad
 I=\int U\,\dd\mu .
\]
The formula for $U$ follows by circle averaging. Compare $\mu$ with the
empirical configuration of the $z_j$, each smeared on a circle of radius
$\varepsilon$. The signed difference has zero total mass. Its logarithmic
energy is nonpositive, by the Gaussian integral representation of the
logarithm and positivity of the Gaussian quadratic form. Expanding the
three energies and bounding the change of a witness potential under
smearing by $\varepsilon/r_i$ yields
\[
 \log|\Delta(z)|^2\le
 2h\sum_j U(z_j)-h^2 I+
 2h^2\varepsilon\sum_i b_i/r_i-h\log\varepsilon .
\]
For repeated nodes the determinant vanishes directly. Take
$\varepsilon=1/n$; the last two terms are $o(N)$. For a lower bound on
$I$, the diagonal interaction contributes $\log r_i$ while the
off-diagonal interaction is at least $\log|u_i-u_j|$.

For completeness, the limiting elementary phase is
\[
 \Phi(x,y)=-2\pi y+\sum_{(a,b,k)} k\left[
 G(x+(b-118)/100,y)-G(x+(a-118)/100,y)-w(\log w-1)\right],
\]
where $w=(b-a)/100$, $G(v,y)=\tfrac v2\log(v^2+y^2)-v+
y\arctan(v/y)$ by continuous extension at $y=0$, and the five triples are
$(0,50,5),(186,236,5),(68,168,-6),(84,152,-4),(86,150,-2)$.
They come from the rising products in the kernel; factorial
normalization gives the width terms. In the upper half-plane the limiting
contour is $x=1.18-c(y)$ with $c$ interpolating the table. The combined
field $F(z)=\Phi(x,y)+(8/5)U(z^2)$ obeys
\[
 F(z)\le P_*=-8443338478/10^9.
\]
Here is how a finite set of inequalities controls the whole contour.
On the 119 panels up to $y=4$, write $z_t=(1-t)z+tw$. Bounds on the
negative second derivative make the field plus a quadratic convex
on each smooth piece. The clipped logarithm is a maximum with a constant;
its switching slopes permit the same chord bound across a switch.
The phase contribution to the chord error is $t(1-t)$ times a sum of
terms $|k||w-z|^2/(2d_v)$ over block ends. Here $d_v$ is a rounded-down
lower bound on distance to the translated origin. A circle contributes
$t(1-t)b_i\big(|w^2-z^2|^2/(2d_i^2)+|w-z|^2/d_i\big)$ to the potential
bound, where $d_i$ is the larger of $r_i$ and a rounded-down minimum
endpoint distance in the square plane. The second term accounts for the
difference between $z_t^2$ and the chord joining the squared endpoints.
Use $t(1-t)\le1/4$, upper endpoint values and outward rounding. The file
\code{ChordBudget.lean} gives the exact integer comparison for each panel.
For $y\ge4$ the vertical ray has $x=327/200$ and
$F(z)\le P_*-6(y-4)$.

Finite-product errors are controlled before taking limits: at scale $D$
one first chooses a positive quadrature parameter to absorb the linear
error into $\delta D$, then takes $n$ large to absorb logarithmic and
constant terms. Explicitly the finite-product allowance is $10D\pi\eta+40/\eta+12(9+\log D)$. The real-coordinate drift is at most $2/D$. The resulting
one-variable integral is at most $\exp(D(P_*+\delta))$ for large $n$
and then sufficiently large truncation height. Taking the $h$-fold
product uses $hD=N$. Stirling bounds give
$\log B_n/N\to 12/5+(8/5)\log(5/8)$. Since $h^2/N=4/5$, the certified
strict inequality
\[
 12/5+(8/5)\log(5/8)+P_*-(4/5)I<-619/100
\]
absorbs the positive margins. This proves the eventual bound \eqref{eq:A}.

\section{Prime costs and the arithmetic certificate}
Let $P=p/n$ and $e_n(p)=\max(0,-\val c_n)$. Bounding
$-\log|c_n|$ by $\sum_{p\le84n} e_n(p)\log p$ reduces the calculation to
three bands. Weighted prime number theorem limits convert piecewise
majorants in $P$ into integrals. Fixed breakpoints, the diagonal prime
$p=n$ where relevant, and higher prime powers are estimated separately.
The normalizing divisor in all three integrals is $2000$.

\paragraph{Large primes.}
For $50<P\le84$ split off the multiples of $p$ in the rational harmonic
part. The regular power-basis part is integral. Singular sites have
multiplicities ten or six and their ranked column capacities decrease by
two per rank. Counting the allowed slots at each level and capping by
$40n$ bounds the determinant valuation. The even basis adds at most
$\max(0,80n-p)$. For $50n<p\le51n$ there are $51n+1-p$ tenfold sites and
$8n$ sixfold sites. At levels 1 and 2 use the rank cap twice. The
remaining slot-level counts per tenfold and sixfold site sum to 42 and
20, respectively: the slot capacities are the positive entries of
$m+4-2a$ for $0\le a<m$, and removing the first two levels subtracts
two from each positive capacity. This gives at most
$240n+42(51n-p)+42$ before the basis correction.
For $51n<p\le59n$ only $59n+1-p$ sixfold sites remain, each with
total capacity 30. Above $59n$ the power-basis matrix is integral.
Consequently $e_n(p)\le n W(p/n)+42$ with the following four
affine majorants:
\[
\begin{array}{c|cccc}
 P&(50,51]&(51,59]&(59,80]&(80,84]\\ \hline
 W(P)&2462-43P&1850-31P&80-P&0 .
\end{array}
\]
Integrating these expressions and dividing by 2000 gives $1671/2000$.

\paragraph{Middle primes.}
For $11/20<P\le50$ and large $n$, $p^2>84n$. Residues fold in opposite
pairs. In fourfold integer coordinates $p'=4\operatorname{num}(P)$,
$d=4\operatorname{den}(P)$ the wall positions are $0,p'/2$ and
$\min(ed\bmod p',\,p'-(ed\bmod p'))$ for $e=16,17,25,34,59$.
Between consecutive walls, evaluate the two residue progressions at a
representative midpoint. Each gives a descending bag of positive slot
weights (at most 16). Multiplicities of sites and aliases come from floor
differences of the original numerator and denominator intervals.
For clarity, at an integer representative $t$ of one progression put
\[
 H(a,b)=\lfloor(bd-t)/p'\rfloor-\lfloor(ad-t)/p'\rfloor,\quad
 m=6+4[\,t<17d\,]+2[\,t<16d\,].
\]
The denominator alias count is
$u=6H(-25,25)+4H(-17,17)+2H(-16,16)-m$ and the numerator count is
$v=5(H(-59,-34)+H(34,59))$. The base valuation is
\[
 v_0=v-u+6\lfloor50d/p'\rfloor+4\lfloor34d/p'\rfloor+
 2\lfloor32d/p'\rfloor-10\lfloor25d/p'\rfloor .
\]
Scan $a=m-1,\ldots,0$. For $1\le e\le m-a$ put $o=m-e-a$ and
$\iota=o$ if $u>0$, otherwise $\iota=\min(o,v)$. With
$l=\lfloor(34d-t)/p'\rfloor$, $r=\lfloor(34d+t)/p'\rfloor$, let $J$
denote $r>0$ when $e$ is odd, and $l\ne r$ otherwise. At each $a$ retain
the running maximum of zero and all scores
$a-v_0+\iota+(e+4)[J]$ considered so far. Its positive values form the
site slots. This rule records the worst valuations allowed by the local
Taylor quotient and finite harmonic sums.
For layer $j$, count all eligible slots across the strips weighted by
their strip width, and cap the population by $40d$. The sum over
$j=1,\ldots,16$ gives the unfolded cost. For a cut $q$ replace a weight
$w_r$ of zero-based rank $r$ by $w_r-2\max(r-q,0)$ and add the row cost
\[
 80d(K-q)-p'\,(K+q+1)(K-q)/2,\qquad K=\lfloor80d/p'\rfloor .
\]
The product in the numerator is even. The minimum of the unfolded
and admissible folded costs, divided by $2000d$, is $V(P)$.
The discrete determinant bound, for the eligible primes off the
exceptional set, is $e_n(p)\le2000nV(P)+123648$. The additive constant
accounts for strip boundaries and disappears after prime summation and
normalization.

On each open certificate interval the ordering of walls, floor counts
and ranks relevant to the expression are stable. Widths are affine in
$P$; adding costs and taking minima, including the layer caps, produces
a concave expression. At the midpoint, choose the active branch
at each minimum and add the selected affine supports at each sum.
Induction on the expression gives an affine majorant that touches it at
the midpoint; integrating this line gives exactly its midpoint value
times the width. Continuity extends the integral estimate to
the closed interval. For a row $[a,b]$ the contribution is at most
$(b-a)V((a+b)/2)$ by the midpoint inequality for concave functions.
The table stores the ceiling of this product times $10^{12}$.
Consecutive rows cover $[11/20,50]$. Their 1,375 ceilings sum to
$5244756877432$, yielding the normalized middle budget.

\paragraph{Small primes.}
For $P\le11/20$ put $\tau=1/P=k+u$ with $0\le u\le1$. Integer
translation leaves the centered residue counts periodic in $u$.
For a precise row formula let $u=n'/d'$ be a row midpoint in lowest
terms. Sort $0,d'$ and the residues $v n'\bmod d'$ for
$v\in\{0,\pm16,\pm17,\pm25,\pm34,\pm59\}$. For consecutive
walls $\alpha,\beta$ define
\[
 T(l,r)=\left\lfloor\frac{2rn'-\alpha-\beta}{2d'}\right\rfloor
 -\left\lfloor\frac{2ln'-\alpha-\beta}{2d'}\right\rfloor
 -\left\lfloor\frac{(r-l)n'}{d'}\right\rfloor .
\]
Take the positive part of $4+6T(-25,25)+4T(-17,17)+2T(-16,16)
-5T(-59,-34)-5T(34,59)$ and call it $v$. The means $F,G$ of
$v,\lfloor(v-1)^2/4\rfloor$, with weights $(\beta-\alpha)/d'$, give
the period scores. Within each of the 706 certified cells the wall
calculation gives the midpoint integral. On a cell $[a,b]$ a majorant
for the normalized contribution is
\[
 \frac{b-a}{50}\left(F S_2(a)+\frac{F+G}{80}S_3(a)\right).
\]
Here $S_j(a)$ sums upward-rounded $(k+a)^{-j}$ from $k=1$ if
$a\ge9/11$ and from $k=2$ otherwise, through $16$, plus the rounded
tail $1/((j-1)(16+a)^{j-1})$. The split at $9/11$ follows from
$\tau\ge20/11$. Truncating first and controlling the very small primes
uniformly justifies the infinite-period limit. Adding the rounded row
costs gives $60166783595/10^{12}$.

The three budgets sum to $6.140423661027$. The PNT permits an arbitrary
positive margin; choosing $10^{-5}$ stays below $6.1405$. Exponentiating
proves \eqref{eq:C}. Combining the two rounded rates gives $-99/2000$.

\section{Formalization, computation and provenance}
The root module \code{OddZetaFive} imports the unconditional statement.
The following module names are relative to \code{OddZetaSix.B40a}.
\begin{center}\small
\begin{tabular}{@{}>{\raggedright\arraybackslash}p{.27\linewidth}>{\raggedright\arraybackslash}p{.67\linewidth}@{}}
\toprule
Part of the argument & Representative modules\\ \midrule
Statement and evaluation & \code{Main}, \code{PolynomialModel}, \code{Identity.Evaluation}\\
Prime specialization & \code{Padic.WindowHermite}, \code{Padic.ConcreteGate}\\
Contour and logarithmic energy & \code{Analysis.VandermondeContour}, \code{Analysis.SignedCircleEnergy}, \code{Analysis.AnalyticRate}\\
Continuous field bounds & \code{Analysis.ChordGeometry}, \code{Analysis.CertifiedField}, \code{Analysis.RayTail}\\
Arithmetic bounds & \code{Arithmetic.MiddleIntegral}, \code{Arithmetic.MiddleBandRate}, \code{Arithmetic.OuterBandsRate}, \code{Arithmetic.ArithRate}\\
Finite data and trust & \code{Certificate.Generated.All},\newline \code{Certificate.ArithmeticNativeAudit}\\
\bottomrule
\end{tabular}
\end{center}
The project pins Lean 4.34.0-rc2 and Mathlib. The local PNT package
selects sources from \emph{PrimeNumberTheoremAnd} (revision
\code{c39a7511}) and \emph{LeanArchitect} (\code{468e8f58}).
Its provenance manifest records the exact revisions, hashes and
source transformations, including a compatibility change in the Wiener
module. Weighted prime limits enter the arithmetic proof via
\code{Arithmetic.WeightedPNT}.

The rebuilt release exports
\begin{flushleft}\small\ttfamily
OddZetaSix.B40a.\\
exists\_irrational\_riemannZeta\_seven\_to\_fifteen
\end{flushleft}
The fresh \texttt{\#print axioms} audit lists \code{propext},
\code{Classical.choice}, \code{Quot.sound}, and one computational axiom:
\begin{flushleft}\small\ttfamily
OddZetaSix.B40a.Cert.Native.middle\_rows\_checked.\\
\_native.native\_decide.ax\_1\_1
\end{flushleft}
The analytic rate, the specialized determinant valuation, the small-prime
rate and the high-prime rate each printed precisely the first three
axioms. The arithmetic rate printed the same four as the top theorem.
The complete output accompanies the source as \code{reports/axioms.log}.
The compiler evaluates 1,375 rational row checks. Their combined Boolean
result enters the middle-band budget and hence the final theorem. A pilot
using one folding witness per row exceeded the two-minute limit on the
first 32-row shard. The retained proof uses native evaluation at this
point. The audit is reproducible by building the package and passing
\code{OddZetaFiveAudit.lean} to Lean. The separate Python verifier uses
integer and rational arithmetic, enclosing logarithms by directed
bounds around correctly rounded standard-library decimal logarithms.
It checks both exact arithmetic row tables, the high-prime integral,
the energy lower bound, each continuous panel budget, and the ray budget.
This corroborates the data, while the Lean computation retains the
trust dependency displayed above.

\begin{thebibliography}{9}
\bibitem{zud} W. Zudilin, \emph{Arithmetic of linear forms involving odd
zeta values}, \url{https://arxiv.org/abs/math/0206176}.
\bibitem{zud2} W. Zudilin, \emph{Irrationality of values of zeta-function},
arXiv:math/0104249 (2001),
\url{https://arxiv.org/abs/math/0104249}.
\bibitem{six} \texttt{janahesse580-star}, \emph{odd-zeta-six-lean},
GitHub revision\newline
\href{https://github.com/janahesse580-star/odd-zeta-six-lean/tree/09201a928982c1e282ba2a1a3156d8f400d4b436}{\code{09201a928982c1e282ba2a1a3156d8f400d4b436}}.
\end{thebibliography}
\end{document}
